% GENERATED FILE. Edit canonical lessons, then run npm run generate:books.
% canonical-source-hash: fnv1a64:30c2d82883272545
% canonical-lessons: MR-C61-sahava, MR-C61-satva, MR-C61-aathva, MR-C61-navva, MR-C61-dahava, MR-R61-sixth-to-tenth

\chapter{The Rest Are Built}
\label{ch:mr-ordinals-sixth-to-tenth}
\hlchaptermodality{61}

\begin{chapteropening}
\textbf{By the end of this chapter:} \emph{I can build any ordinal from sixth upwards by putting -vaa on the number, and I know the one place between five and ten where the number underneath is an older form.}

Complete MR-R61-sixth-to-tenth and run pahilaa to dahaavaa, then name the one of the ten that bent the rule and how.
\end{chapteropening}

\section[sahāvā]{\mr{सहावा} (sahāvā) — sixth}

\label{lesson:MR-C61-sahava}



\begin{quote}
Four recalls, at four distances.

\begin{itemize}
\raggedright
  \item \emph{Recall:} say \emph{pāchvā} — \textbf{R1}, two lessons back
  \item \emph{Recall:} say \emph{tisrā} — \textbf{R2}, five lessons back
  \item \emph{Recall:} say \emph{terā} — \textbf{R3}, twenty lessons back
  \item \emph{Recall:} say \emph{var} — \textbf{R4}, eighty lessons back
\end{itemize}
\end{quote}

\subsection*{You'll want to know: \mr{सहावा}}
\begin{quote}
\textbf{\mr{सहावा}} — \emph{sahāvā} — \textbf{sixth}
\end{quote}

Before you read the word, work it out. You have \textbf{\mr{सहा}} and you have \textbf{-\mr{वा}}.

\textbf{\mr{सहा}} + \textbf{\mr{वा}} = \textbf{\mr{सहावा}}. Nothing was worn down. Nothing was replaced.

This is the first time in this book that a new word has been handed to you already known. It will happen four more times in this chapter.

\begin{grammarlens}[title={Grammar Lens: what a rule buys}]
\begin{quote}
\textbf{\mr{सहावं} \mr{घर}} — \emph{sahāvaṁ ghar} — \textbf{the sixth house}
\end{quote}

Chapter sixty cost you four separate words to reach \emph{fourth}. This chapter costs you one ending, already paid for, to reach \emph{tenth} — and then \emph{twentieth}, and \emph{hundredth}, and every ordinal Marathi will ever need.

That is what the difference between an inherited word and a built one is worth, and it is why the chapter before this one taught its four the hard way rather than pretending they followed a pattern.
\end{grammarlens}

\subsection*{Your turn}
Say these aloud:

\begin{itemize}
\raggedright
  \item \emph{sahā}, then \emph{sahāvā}
  \item \emph{pāchvā}, \emph{sahāvā} — fifth, sixth
  \item what changed inside the number when the ending went on — \textbf{nothing}
\end{itemize}

\begin{tcolorbox}[breakable,colback=teal!4,colframe=teal!35!black,title={Before you move on}]
What is \textbf{\mr{सहा}} plus \textbf{-\mr{वा}}? (\textbf{\mr{सहावा}}.) Did you have to be told this word? (\textbf{No.})

Sources: \href{https://www.omniglot.com/language/numbers/marathi.htm}{Omniglot: Numbers in Marathi}.
\end{tcolorbox}
\section[sātvā]{\mr{सातवा} (sātvā) — seventh}

\label{lesson:MR-C61-satva}



\begin{quote}
Four recalls, at four distances.

\begin{itemize}
\raggedright
  \item \emph{Recall:} say \emph{sahāvā} — \textbf{R1}, one lesson back
  \item \emph{Recall:} say \emph{pahilā} — \textbf{R2}, five lessons back
  \item \emph{Recall:} say \emph{chaudā} — \textbf{R3}, twenty lessons back
  \item \emph{Recall:} say \emph{khālī} — \textbf{R4}, eighty lessons back
\end{itemize}
\end{quote}

\subsection*{You'll want to know: \mr{सातवा}}
\begin{quote}
\textbf{\mr{सातवा}} — \emph{sātvā} — \textbf{seventh}
\end{quote}

\textbf{\mr{सात}} + \textbf{\mr{वा}}. Say it before you read it.

One small thing in the sound. \textbf{\mr{सहा}} ends in a vowel, so \textbf{\mr{सहावा}} comes out in three beats — \emph{sa-hā-vā}. \textbf{\mr{सात}} ends in a consonant, so \textbf{\mr{सातवा}} comes out in two — \emph{sāt-vā}. The ending did not change; the number it landed on did.

\begin{grammarlens}[title={Grammar Lens: the dictionaries agree with you}]
Look up \textbf{\mr{पहिला}} in Molesworth's 1857 dictionary and it is there, with its own definition. Look up \textbf{\mr{तिसरा}}, and \textbf{\mr{चौथा}} — both there, each with its own entry.

Look up \textbf{\mr{सातवा}} and there is nothing. The dictionary has no such word.

That is not an oversight. A dictionary records what a reader cannot work out, and by 1857 nobody thought \textbf{\mr{सातवा}} needed recording, because anybody with \textbf{\mr{सात}} and the ending already had it. \textbf{The seam you were shown last chapter is visible in the dictionary's own table of contents.}
\end{grammarlens}

\subsection*{Your turn}
Say these aloud:

\begin{itemize}
\raggedright
  \item \emph{sahāvā}, \emph{sātvā}
  \item \emph{sāt} and then \emph{sātvā}
  \item which four ordinals a dictionary DOES print — \emph{pahilā dusrā tisrā chauthā}
\end{itemize}

\begin{tcolorbox}[breakable,colback=teal!4,colframe=teal!35!black,title={Before you move on}]
What is \textbf{\mr{सात}} plus \textbf{-\mr{वा}}? (\textbf{\mr{सातवा}}.) Why is that word missing from the dictionary? (\textbf{Because it can be worked out.})

Sources: \href{https://dsal.uchicago.edu/cgi-bin/app/molesworth\_query.py?qs=\%E0\%A4\%B8\%E0\%A4\%BE\%E0\%A4\%A4\%E0\%A4\%B5\%E0\%A4\%BE\&searchhws=yes}{Molesworth, \emph{A Dictionary, Marathi and English} — search for \mr{सातवा} (DSAL)}; \href{https://www.omniglot.com/language/numbers/marathi.htm}{Omniglot: Numbers in Marathi}.
\end{tcolorbox}
\section[āṭhvā]{\mr{आठवा} (āṭhvā) — eighth}

\label{lesson:MR-C61-aathva}



\begin{quote}
Four recalls, at four distances.

\begin{itemize}
\raggedright
  \item \emph{Recall:} say \emph{sātvā} — \textbf{R1}, one lesson back
  \item \emph{Recall:} say \emph{chauthā} — \textbf{R2}, five lessons back
  \item \emph{Recall:} say \emph{pandhrā} — \textbf{R3}, twenty lessons back
  \item \emph{Recall:} say \emph{khoolī} — \textbf{R4}, eighty-six lessons back
\end{itemize}
\end{quote}

\subsection*{You'll want to know: \mr{आठवा}}
\begin{quote}
\textbf{\mr{आठवा}} — \emph{āṭhvā} — \textbf{eighth}
\end{quote}

\begin{quote}
\textbf{\mr{आठवी} \mr{खोली}} — \emph{āṭhvī khoolī} — \textbf{the eighth room}
\end{quote}

\textbf{\mr{आठ}} + \textbf{\mr{वा}}, with the retroflex \textbf{\mr{ठ}} carried through untouched.

\begin{grammarlens}[title={Grammar Lens: the pair to keep apart}]
Two words in this set both end in a \textbf{-\mr{था}}/\textbf{-\mr{वा}} you might blur together:

\noindent
\begin{tabularx}{\linewidth}{@{}>{\raggedright\arraybackslash}X>{\raggedright\arraybackslash}X@{}}
\toprule
\textbf{} & \textbf{} \\
\midrule
\textbf{\mr{चौथा}} (\emph{chauthā}) & fourth — an inherited word, \textbf{-\mr{था}} \\
\textbf{\mr{आठवा}} (\emph{āṭhvā}) & eighth — a built word, \textbf{\mr{आठ}} + \textbf{-\mr{वा}} \\
\bottomrule
\end{tabularx}

They are not the same ending, and the test is quick: take the ending off and see whether a number is left standing. \textbf{\mr{आठवा}} minus \textbf{-\mr{वा}} leaves \textbf{\mr{आठ}}. \textbf{\mr{चौथा}} minus \textbf{-\mr{था}} leaves \textbf{\mr{चौ}}, which is not a Marathi number.
\end{grammarlens}

\subsection*{Your turn}
Say these aloud:

\begin{itemize}
\raggedright
  \item \emph{āṭh}, then \emph{āṭhvā}
  \item \emph{sahāvā sātvā āṭhvā} — sixth, seventh, eighth
  \item \emph{āṭhvī khoolī} — the eighth room
\end{itemize}

\begin{tcolorbox}[breakable,colback=teal!4,colframe=teal!35!black,title={Before you move on}]
What is \textbf{\mr{आठ}} plus \textbf{-\mr{वा}}? (\textbf{\mr{आठवा}}.) Is \textbf{\mr{चौथा}} built the same way? (\textbf{No} — take \textbf{-\mr{था}} off and no number is left.)

Sources: \href{https://www.omniglot.com/language/numbers/marathi.htm}{Omniglot: Numbers in Marathi}.
\end{tcolorbox}
\section[navvā]{\mr{नववा} (navvā) — ninth}

\label{lesson:MR-C61-navva}



\begin{quote}
Four recalls, at four distances.

\begin{itemize}
\raggedright
  \item \emph{Recall:} say \emph{āṭhvā} — \textbf{R1}, one lesson back
  \item \emph{Recall:} say \emph{pāchvā} and the ending that made it — \textbf{R2}, five lessons back
  \item \emph{Recall:} say \emph{soḷā} — \textbf{R3}, twenty lessons back
  \item \emph{Recall:} say \emph{javaḷ} — \textbf{R4}, eighty lessons back
\end{itemize}
\end{quote}

\subsection*{You'll want to know: \mr{नववा}}
\begin{quote}
\textbf{\mr{नववा}} — \emph{navvā} — \textbf{ninth}
\end{quote}

Work it out first, the way the last three lessons let you. Nine is \textbf{\mr{नऊ}} (\emph{naū}), so the ending should give \textbf{\mr{नऊवा}}.

It does not. The word is \textbf{\mr{नववा}} — \emph{nav-vā}.

\begin{grammarlens}[title={Grammar Lens: one tooth missing, and where it went}]
\noindent
\begin{tabularx}{\linewidth}{@{}>{\raggedright\arraybackslash}X>{\raggedright\arraybackslash}X>{\raggedright\arraybackslash}X@{}}
\toprule
\textbf{number} & \textbf{expected} & \textbf{actual} \\
\midrule
\textbf{\mr{सहा}} & \textbf{\mr{सहावा}} & \textbf{\mr{सहावा}} \\
\textbf{\mr{सात}} & \textbf{\mr{सातवा}} & \textbf{\mr{सातवा}} \\
\textbf{\mr{आठ}} & \textbf{\mr{आठवा}} & \textbf{\mr{आठवा}} \\
\textbf{\mr{नऊ}} & \textbf{\mr{नऊवा}} & \textbf{\mr{नववा}} \\
\bottomrule
\end{tabularx}

The rule held four times and slipped once. What slipped is not the ending — it is still \textbf{-\mr{वा}} — but the shape of the number underneath it. Wiktionary traces \textbf{\mr{नऊ}} back through Prakrit \emph{nava} to Sanskrit \textbf{\mr{नव}}, and it is that older \textbf{\mr{नव}} that the ordinal is standing on. \textbf{The ending froze the number as it was when the ordinal was made}, and the cardinal went on changing without it.

So what you can say now is a little longer than what last chapter said. From five up, an ordinal is the number with \textbf{-\mr{वा}} behind it — and at nine, the number that turns up underneath is an older shape of itself. One adjustment in the whole system, worth carrying as one adjustment rather than as a second pattern.
\end{grammarlens}

\subsection*{Your turn}
Say these aloud:

\begin{itemize}
\raggedright
  \item \emph{sahāvā sātvā āṭhvā navvā} — sixth to ninth
  \item which one of the four did not come out as expected — \emph{\textbf{navvā}}
  \item the part that DID stay the same in all four — \emph{\textbf{-vā}}
\end{itemize}

\begin{tcolorbox}[breakable,colback=teal!4,colframe=teal!35!black,title={Before you move on}]
What is ninth? (\textbf{\mr{नववा}}.) Which part of the rule bent, the ending or the number? (\textbf{The number} — the ending held.)

Sources: \href{https://www.omniglot.com/language/numbers/marathi.htm}{Omniglot: Numbers in Marathi}; \href{https://en.wiktionary.org/wiki/\%E0\%A4\%A8\%E0\%A4\%8A}{Wiktionary: \mr{नऊ}}.
\end{tcolorbox}
\section[dahāvā]{\mr{दहावा} (dahāvā) — tenth}

\label{lesson:MR-C61-dahava}



\begin{quote}
Four recalls, at four distances.

\begin{itemize}
\raggedright
  \item \emph{Recall:} say \emph{navvā} — \textbf{R1}, one lesson back
  \item \emph{Recall:} say \emph{chauthā} — \textbf{R2}, seven lessons back
  \item \emph{Recall:} say \emph{bārā} — \textbf{R3}, twenty-seven lessons back
  \item \emph{Recall:} say \emph{samor} — \textbf{R4}, eighty lessons back
\end{itemize}
\end{quote}

\subsection*{You'll want to know: \mr{दहावा}}
\begin{quote}
\textbf{\mr{दहावा}} — \emph{dahāvā} — \textbf{tenth}
\end{quote}

\textbf{\mr{दहा}} + \textbf{\mr{वा}}, with the number whole and the ending behind it, exactly as sixth, seventh and eighth were.

\begin{grammarlens}[title={Grammar Lens: this is where the set stops, and the ending does not}]
\emph{Tenth} is where this book's ordinal list ends. It is \textbf{not} where Marathi's ends.

You already count to twenty, and every one of those numbers takes the same ending. Here is one you can build without help:

\begin{quote}
\textbf{\mr{बारा}} $\to$ \textbf{\mr{बारावा}} — \emph{bārāvā}, twelfth
\end{quote}

That word is not on this book's list and never will be. You can say it anyway, which is the whole difference between a list and a rule.

\textbf{What you own after this chapter is not ten words. It is four words and an ending} — and the ending goes on working after the book stops listing.
\end{grammarlens}

\subsection*{Your turn}
Say these aloud:

\begin{itemize}
\raggedright
  \item \emph{dahā}, then \emph{dahāvā}
  \item \emph{sahāvā sātvā āṭhvā navvā dahāvā} — sixth to tenth
  \item \emph{bārāvā} — twelfth, a word this book never taught you
\end{itemize}

\begin{tcolorbox}[breakable,colback=teal!4,colframe=teal!35!black,title={Before you move on}]
What is \textbf{\mr{दहा}} plus \textbf{-\mr{वा}}? (\textbf{\mr{दहावा}}.) Does the ending stop at ten? (\textbf{No} — the book's list does; the ending does not.)

Sources: \href{https://www.omniglot.com/language/numbers/marathi.htm}{Omniglot: Numbers in Marathi}.
\end{tcolorbox}
\section[sahāvā · sātvā · āṭhvā · navvā · dahāvā]{Sixth to tenth, and what the dictionary knew}

\label{lesson:MR-R61-sixth-to-tenth}



\begin{quote}
Four recalls, at four distances, before the run.

\begin{itemize}
\raggedright
  \item \emph{Recall:} say \emph{dahāvā} — \textbf{R1}, one lesson back
  \item \emph{Recall:} say \emph{sahāvā} — \textbf{R2}, five lessons back
  \item \emph{Recall:} say \emph{satrā} — \textbf{R3}, twenty lessons back
  \item \emph{Recall:} say \emph{lāmb} — \textbf{R4}, eighty lessons back
\end{itemize}
\end{quote}

\subsection*{Your turn}
Say these aloud:

\begin{itemize}
\raggedright
  \item \emph{pahilā} … \emph{dahāvā} — the whole ten
  \item the four that had to be memorised — \emph{pahilā dusrā tisrā chauthā}
  \item the six that did not — \emph{pāchvā sahāvā sātvā āṭhvā navvā dahāvā}
\end{itemize}

\begin{grammarlens}[title={What you've built across these two chapters}]
Ten ordinals. \textbf{Four were words and six were arithmetic} — and of the six, five came out exactly as the rule predicted and one, \textbf{\mr{नववा}}, kept an older shape of its number.

There is a way to check that split without taking this book's word for it. Open Molesworth's \emph{Dictionary, Marathi and English} — the 1857 one this chapter has been citing — and look up all ten:

\noindent
\begin{tabularx}{\linewidth}{@{}>{\raggedright\arraybackslash}X>{\raggedright\arraybackslash}X@{}}
\toprule
\textbf{looked up} & \textbf{in the dictionary?} \\
\midrule
\textbf{\mr{पहिला}} & yes, p. 497 \\
\textbf{\mr{दुसरा}} & yes, p. 420 \\
\textbf{\mr{तिसरा}} & yes, p. 381 \\
\textbf{\mr{चौथा}} & yes, p. 296 \\
\textbf{\mr{सहावा}} · \textbf{\mr{सातवा}} · \textbf{\mr{आठवा}} · \textbf{\mr{नववा}} · \textbf{\mr{दहावा}} & \textbf{no entry, any of them} \\
\textbf{\mr{पाचवा}} & an entry, but for a different word entirely — a returning sickness \\
\bottomrule
\end{tabularx}

Four in, the built ones out. A dictionary is a list of what a reader cannot work out for themselves, so the seam this book drew between \textbf{\mr{चौथा}} and \textbf{\mr{पाचवा}} was drawn in Marathi lexicography a century and a half before it was drawn here.
\end{grammarlens}

\begin{tcolorbox}[breakable,colback=teal!4,colframe=teal!35!black,title={Before you move on}]
Which four ordinals does an 1857 dictionary print, and why exactly those four? (\textbf{\mr{पहिला}, \mr{दुसरा}, \mr{तिसरा}, \mr{चौथा} — because they are the four that cannot be predicted.})
\end{tcolorbox}
